% !TEX encoding   = UTF8
% !TEX spellcheck = ru_RU

%%==================
\Chapter{Заключение}
%%==================

%% Согласно ГОСТ Р 7.0.11-2011:
%% 5.3.3 В заключении диссертации излагают итоги выполненного исследования, рекомендации, перспективы дальнейшей разработки темы.
%% 9.2.3 В заключении автореферата диссертации излагают итоги данного исследования, рекомендации и перспективы дальнейшей разработки темы.
%% Поэтому имеет смысл сделать эту часть общей и загрузить из одного файла в автореферат и в диссертацию:

Основным результатом работы является создание приложения для быстрого расчета аэродинамических характеристик профиля по его геометрическим параметрам.

Для достижения поставленной цели были решены 4 задачи, по которым можно сделать следующие выводы:
\begin{enumerate}
  \item Изучено не менее 10 источников, посвященных методам параметризации геометрии. В рассмотренных источниках отмечается, что метод CST удовлетворяет основным требованиям к методу параметризации геометрии: гибкости, экономичности и устойчивости.
  \item Решена задача аппроксимации профиля, заданного точками. Показано, что задача сводится к системе линейных уравнений, которая может быть решена методом наименьших квадратов.
  \item Для обучения суррогатной модели, с помощью CFD расчетов была сгенерирована обучающая выборка. Для формирования независимых признаков и области их варьирования используется метод главных компонент и база данных профилей UIUC.
  \item В качестве суррогатной модели опробованы полиномиальные модели 3, 4 и 5-ой степеней, а также модель кригинга. Выбор вида суррогатной модели производился методом скользящего контроля.
  Лучшей себя показала модель кригинга, получив RMSE (корень среднеквадратичной ошибки) равную 0.6 каунтов по $C_x$ и 1.5 каунта по $C_y$.
\end{enumerate}

Естественным продолжением данной работы является обобщение решенной задачи на несимметричные профили. Помимо этого, планируется расширить возможности суррогатной модели, а именно, добавить параметры потока в качестве ее варьируемых переменных. Хорошим дополнением к данной работе будет задача о предсказании распределенных аэродинамических характеристик, например, распределения коэффициента давления на поверхности профиля.
